% !TeX root = ../../Book2.tex
\chapter{椭圆正则性}
\label{b2:ch:16}
\BookChapterNumber{b2:ch:16}
% 本章摘要（不计入编号）
\begin{chapterabstract}
本章研究能量弱解何时具有更高正则性。差商方法在不预设二阶导数的前提下，建立 Lipschitz 系数方程的内部 \(H^2\) 估计；逐项控制求导误差后，得到高阶 Sobolev 正则性和光滑数据的内部光滑性。边界部分通过平坦化、切向差商和方程恢复法向导数，证明 \(C^{1,1}\) 区域上的 Dirichlet 二阶估计。最后仅用有界可测系数建立 Caccioppoli 不等式及截断迭代的局部有界性，并说明 Harnack 比较的完整线性证明如何在附录 H 补齐，随后推出 Hölder 连续性。
\end{chapterabstract}

% 本章目标（不计入编号）
\begin{learninggoals}
\begin{itemize}
\item\label{b2:item:16-goal-difference}使用离散分部积分和一致差商界构造二阶弱导数，指出系数导数的作用。
\item\label{b2:item:16-goal-bootstrap}逐步核验微分后右端的空间，结合 Sobolev 嵌入判断经典正则性。
\item\label{b2:item:16-goal-boundary}理解边界平坦化、切向估计及纯法向恢复，并用凹角反例检验全局结论。
\item\label{b2:item:16-goal-low}在可测系数下证明局部能量界和有界性，区分本章已证结论与附录 H 补齐的 Harnack 比较证明。
\end{itemize}
\end{learninggoals}

上一章的弱解可以没有经典二阶导数，但方程并非只是一条约束：它还把不同方向和不同尺度的变化联系起来。若系数变化受控，差商误差可以吸收，解便得到新的弱导数；若系数只是可测，导数提升可能失败，标量解的高值集合却仍受能量限制。本章分别建立这两类估计，避免把“椭圆正则性”当作无条件的光滑性口号。

\ProseParagraphBreak

本章普通主线承接第十五章的测试等式，使用第一卷第二十四章的 Sobolev 与 Poincaré 不等式。系统参考可读 Lawrence C. Evans 的《Partial Differential Equations》\cite{Evans2010Partial}以及 David Gilbarg 与 Neil S. Trudinger 的《Elliptic Partial Differential Equations of Second Order》\cite{Gilbarg2001Elliptic}。下面对差商、边界 \(H^2\) 与局部有界性给出完整证明；末节先给可测系数 Harnack 比较的证明概要，完整的线性比较证明及其所需指数尾见附录 H、C；不把此处略去的步骤隐藏在“标准估计”中。

% !TeX root = ../../Book2.tex
\section{差商方法}
\label{b2:sec:1601}

能量解只有一阶弱导数，不能直接对方程再求导。差商以两个平移值之差代替导数，先在已有的函数空间中操作，再由一致估计恢复二阶弱导数。这样既避免预设待证正则性，也清楚显示系数的变化会在估计中产生什么误差。

\subsection{差商}
\label{b2:subsec:160101}

设整数 \(d\geq1\)，\(\Omega\subset\mathbb R^d\) 为开集，\(u\in L^1_{\mathrm{loc}}(\Omega;\mathbb K)\)，其中 \(\mathbb K=\mathbb R\) 或 \(\mathbb C\)，\(k\in\{1,\ldots,d\}\)、\(h\in\mathbb R\setminus\{0\}\)。固定坐标方向 \(e_k\)，定义 \term[chashang]{差商}[difference quotient]
\[D_h^ku(x)=\frac{u(x+he_k)-u(x)}h\qquad(h\ne0).\]
只在所需平移仍留在定义域内时使用这一表达式。若 \(f,g\in C_c^\infty(\Omega)\)，并且它们支集的所需 \(h\) 平移都仍在 \(\Omega\) 内，换元给出离散分部积分；该式也适用于满足这些支撑条件的 \(L^2\) 函数
\[\int fD_{-h}^kg=-\int(D_h^kf)g.\]
其中 \(D_{-h}^kg=(g(x-he_k)-g(x))/(-h)\)；负号和负步长必须同时处理。乘积公式是
\[D_h^k(fg)=f(x+he_k)D_h^kg+(D_h^kf)g(x).\]

\Needspace{6\baselineskip}
\begin{lemma}[差商与弱导数]
\label{b2:lem:16-difference-quotient}
设整数 \(d\geq1\)，\(U\Subset V\subset\mathbb R^d\) 为非空开集，标量域 \(\mathbb K=\mathbb R\) 或 \(\mathbb C\)，\(k\in\{1,\ldots,d\}\)，\(h\in\mathbb R\setminus\{0\}\)，\(D_h^ku(x)=[u(x+he_k)-u(x)]/h\)，\(|h|<\operatorname{dist}(U,\partial V)\)（\(V=\mathbb R^d\) 时该距离取 \(+\infty\)）。若 \(u\in H^1(V;\mathbb K)\)，则
\[\lVert D_h^ku\rVert_{L^2(U)}\le\lVert\partial_ku\rVert_{L^2(V)}.\]
反之，若 \(u\in L^2(V;\mathbb K)\)，且存在 \(0\leq C<\infty\)、\(h_0>0\)，使 \(\lVert D_h^ku\rVert_{L^2(U)}\le C\) 对所有 \(0<|h|<h_0\) 且 \(|h|<\operatorname{dist}(U,\partial V)\) 成立，则 \(u\) 在 \(U\) 中有弱导数 \(\partial_ku\in L^2(U;\mathbb K)\)，其范数不超过 \(C\)。
\end{lemma}
\begin{proof}
对光滑 \(u\)，有 \(D_h^ku(x)=\int_0^1\partial_ku(x+the_k)\,\mathrm dt\)。Jensen 不等式、积分换序和各平移集合包含于 \(V\) 给出所列界。一般 \(H^1\) 函数先在包住所有相关线段的内部区域磨光，再以 \(H^1\) 收敛取极限，结论保持。

证明反向时，取 \(h_n\to0\)。\(L^2(U)\) 中的有界差商有弱收敛子列，记极限为 \(g\)。对 \(\varphi\in C_c^\infty(U)\)，离散分部积分和测试函数差商的一致收敛给出
\[\int_Ug\varphi=\lim_n\int_U(D_{h_n}^ku)\varphi
 =-\lim_n\int_VuD_{-h_n}^k\varphi=-\int_Uu\partial_k\varphi.\]
故 \(g\) 正是弱导数，弱下半连续性给出范数界。弱导数的唯一性说明所获对象不依赖抽取的子列。
\end{proof}

引理中的第一项是已有导数给出的上界，第二项则把统一差商界变成导数存在性；正则性证明需要把这两个方向配合使用。仅对某一个固定步长估计当然不能推出可微性。还应保留内部集合 \(U\Subset V\)：在边界附近任意平移会跑出区域，边界估计因此需要新的处理。

\subsection{内部 \texorpdfstring{\(H^2\)}{H2} 估计}
\label{b2:subsec:160102}

\Needspace{6\baselineskip}
\begin{theorem}[内部二阶估计]
\label{b2:thm:16-interior-h2}
设整数 \(d\geq1\)，\(\Omega\subset\mathbb R^d\) 为非空开集，\(r>0\)，\(B_{2r}=B(x_0,2r)\Subset\Omega\)，实矩阵场 \(A\in W^{1,\infty}(B_{2r};\mathbb R^{d\times d})\) 满足 \(\lVert A\rVert_\infty\le\Lambda<\infty\)，且对几乎处处 \(x\in B_{2r}\)、全部 \(\xi\in\mathbb R^d\) 有 \(\xi^{\mathsf T}A(x)\xi\ge\lambda|\xi|^2\)，其中 \(\lambda>0\)，矩阵取欧氏算子范数。若实值 \(u\in H^1(B_{2r};\mathbb R)\)、\(f\in L^2(B_{2r};\mathbb R)\) 且
\[\int_{B_{2r}}A\nabla u\cdot\nabla\varphi=\int_{B_{2r}}f\varphi
 \qquad(\varphi\in C_c^\infty(B_{2r};\mathbb R)),\]
则 \(u\in H^2(B_r)\)，并有
\[\lVert D^2u\rVert_{L^2(B_r)}
 \le C\left(\lVert f\rVert_{L^2(B_{2r})}
 +(r^{-1}+\lVert\nabla A\rVert_\infty)\cdot\lVert\nabla u\rVert_{L^2(B_{2r})}\right),\]
其中 \(C\) 只依赖 \(d,\lambda,\Lambda\)。
\end{theorem}
\begin{proof}
取光滑截断 \(0\le\eta\le1\)，使其在 \(B_r\) 上为 \(1\)，支集包含于 \(B_{3r/2}\)，且 \(|\nabla\eta|\le C_d/r\)。固定方向 \(k\)，取 \(|h|<r/8\)，记 \(w=D_h^ku\)。对这个固定步长，\(w\) 属于较小球上的 \(H^1\)，所以 \(v=-D_{-h}^k(\eta^2w)\) 属于 \(H_0^1(B_{2r})\)。型与右端均连续，密度保证可以用 \(v\) 测试。

将差商移到通量，得到
\[
 \int D_h^k(A\nabla u)\cdot\nabla(\eta^2w)
 =-\int fD_{-h}^k(\eta^2w).
\]
乘积公式给出 \(D_h^k(A\nabla u)=A(x+he_k)\nabla w+(D_h^kA)\nabla u\)。其中主项为 \(\int\eta^2A(x+he_k)\nabla w\cdot\nabla w\)，下界是 \(\lambda X^2\)，记 \(X=\lVert\eta\nabla w\rVert_2\)。设 \(U=\lVert\nabla u\rVert_{L^2(B_{2r})}\)、\(F=\lVert f\rVert_{L^2(B_{2r})}\)、\(K=\lVert\nabla A\rVert_\infty\)。引理给出 \(\lVert w\rVert_{L^2(B_{7r/4})}\le U\)，且 Lipschitz 控制给出 \(\lVert D_h^kA\rVert_\infty\le K\)。

展开 \(\nabla(\eta^2w)=\eta^2\nabla w+2\eta w\nabla\eta\)。其余主部交叉项、系数差商项及右端分别被
\[C\Lambda r^{-1}UX,\qquad KUX+CKr^{-1}U^2,
 \qquad CF(X+r^{-1}U)\]
控制。最后一界再次使用差商引理于紧支集函数 \(\eta^2w\)，即 \(\lVert D_{-h}^k(\eta^2w)\rVert_2\le\lVert\partial_k(\eta^2w)\rVert_2\)。逐项用 Young 不等式吸收不超过 \(\lambda X^2/2\) 的部分，余项得到
\[X^2\le C\bigl(F^2+(r^{-1}+K)^2U^2\bigr),\]
常数与 \(h\) 无关。因为 \(\eta=1\) 于 \(B_r\)，各 \(D_h^k\partial_j u\) 在该球的 \(L^2\) 范数一致有界。应用引理的反向得到全部 \(\partial_k\partial_j u\in L^2(B_r)\)，并对有限个指标求和即得估计。
\end{proof}

这个证明没有先假定二阶导数存在。测试函数中的差商仍只要求 \(u\in H^1\)；真正的新导数是在获得与步长无关的界之后才出现。系数的 Lipschitz 范数也有明确来源：它控制 \(D_hA\)。若系数只有本性上界，这个误差可能按 \(|h|^{-1}\) 增大，证明就不能闭合。

这里的矩阵系数取实值，解和右端的复值版本可直接分开处理：若 \(u=u_1+\mathrm i u_2\)、\(f=f_1+\mathrm i f_2\)，同一实矩阵 \(A\) 使两个实方程分别成立。各用一次定理，再把两个平方估计相加，就得到同一形式、仍只依赖 \(d,\lambda,\Lambda\) 的复值估计。下节的高阶提升也可逐次按这个办法处理复值数据；一般复矩阵的椭圆条件与配对须另作规定。

\subsection{局部正则性}
\label{b2:subsec:160103}

设整数 \(d\geq1\)，\(\Omega\subset\mathbb R^d\) 为非空开集，\(A\in W_{\mathrm{loc}}^{1,\infty}(\Omega;\mathbb R^{d\times d})\)，实值 \(f\in L_{\mathrm{loc}}^2(\Omega)\)、\(u\in H_{\mathrm{loc}}^1(\Omega)\) 满足 \(-\operatorname{div}(A\nabla u)=f\) 的局部测试恒等式。给定非空开集 \(U\Subset V\Subset\Omega\)，假设对几乎处处 \(x\in V\)、全部 \(\xi\in\mathbb R^d\) 有 \(\xi^{\mathsf T}A(x)\xi\geq\lambda_V|\xi|^2\)，其中 \(\lambda_V>0\)。用等半径内部球覆盖 \(\overline U\)，使加倍球均在 \(V\) 内、重叠次数只由 \(d\) 控制，再将球上估计平方求和，得到
\[\lVert u\rVert_{H^2(U)}\le C_{U,V,A}\bigl(\lVert f\rVert_{L^2(V)}+\lVert u\rVert_{H^1(V)}\bigr).\]
常数可只按 \(d\)、\(\operatorname{dist}(U,\partial V)^{-1}\)、\(\lambda_V\) 及 \(\|A\|_{W^{1,\infty}(V)}\) 选取。若每个紧含子域都有这样的椭圆下界，则局部 Lipschitz 系数和局部平方可积右端把能量解提高到 \(H_{\mathrm{loc}}^2\)。常数在区域逼近边界时可能发散，不能由此直接得到全域 \(H^2\) 估计。

\ProseParagraphBreak

若方程还含 \(b\cdot\nabla u+cu\)，且低阶系数有界，则把它移到右端：\(\widetilde f=f-b\cdot\nabla u-cu\in L_{\mathrm{loc}}^2\)。内部二阶估计仍适用，右端多出已有的 \(H^1\) 范数。这一步不需要对 \(b,c\) 求导；只有继续提高阶数时才会用到它们更高的正则性。

\ProseParagraphBreak

一个对照例是前章的分段导热系数：弱解通量连续而梯度跳跃，故二阶分布导数含点质量，不在 \(L^2\)。它说明系数正则性要求具有真实数学内容。另一方面，对常系数 Laplace 方程，本定理不要求 \(f\) 连续，只要求 \(f\in L^2\)；所得结论是二阶弱导数平方可积，而不是二阶经典导数连续。正则性结论应始终与数据所在的空间一起读取。

% !TeX root = ../../Book2.tex
\section{自举法}
\label{b2:sec:1602}

得到二阶弱导数后，方程可以合法地再作分布微分。但每次微分都会同时作用于系数，产生新的右端。自举的关键不是一句“反复求导”，而是证明这些新项具有足够的可积性，使上一节的估计能够再次使用。

\subsection{系数正则性}
\label{b2:subsec:160201}

以 \(-\operatorname{div}(A\nabla u)=f\) 为例，形式求 \(x_k\) 导数得到
\[
 -\operatorname{div}(A\nabla\partial_ku)
 =\partial_kf+\operatorname{div}((\partial_kA)\nabla u).
\]
若只有 \(A\in W^{1,\infty}\)、\(f\in L^2\)，右端一般只属于 \(H^{-1}\)，不能再次套用要求 \(L^2\) 右端的二阶定理。若进一步 \(A\in W^{2,\infty}\)、\(f\in H^1\)，并且已经知道 \(u\in H^2\)，则 \((\partial_kA)\nabla u\in H^1\)，其散度属于 \(L^2\)，因而可以把 \(\partial_ku\) 提高到 \(H^2\)。

\ProseParagraphBreak

系数的这一级额外导数不可随意省略。例如在一维取正的 Lipschitz 系数 \(a(x)=1+|x|\)，令 \(u'(x)=1/a(x)\)。则 \(-(au')'=0\)，\(u\in H_{\mathrm{loc}}^2\)，但 \(u''\) 在零点有跳跃，故 \(u\notin H_{\mathrm{loc}}^3\)。右端已光滑，解仍受到系数不光滑的限制。这个例子精确对应上一行中对 \(\partial_kA\) 再求导的需要。

\ProseParagraphBreak

本节所列 \(W^{m+1,\infty}\) 条件是便于统一证明的一组充分条件，并非每个具体方程的最弱假设。常系数的导数项全部消失；若某些系数仅沿特定方向变化，也可能先在其他方向取得更高正则性。使用定理时可以利用这些特殊结构，但需要重做相应的乘积估计，不能把一般结论的假设凭经验删除。

\subsection{高阶 Sobolev 正则性}
\label{b2:subsec:160202}

\begin{theorem}[内部高阶正则性]
\label{b2:thm:16-higher-regularity}
设 \(m\) 为非负整数，\(A\in W_{\mathrm{loc}}^{m+1,\infty}(\Omega;\mathbb R^{d\times d})\) 局部有界且一致椭圆，\(f\in H_{\mathrm{loc}}^m(\Omega)\)。若 \(u\in H_{\mathrm{loc}}^1(\Omega)\) 弱解 \(-\operatorname{div}(A\nabla u)=f\)，则 \(u\in H_{\mathrm{loc}}^{m+2}(\Omega)\)。对 \(U\Subset V\Subset\Omega\)，有
\[\lVert u\rVert_{H^{m+2}(U)}
 \le C\bigl(\lVert f\rVert_{H^m(V)}+\lVert u\rVert_{H^1(V)}\bigr),\]
其中 \(C\) 依赖区域间距、椭圆下界和 \(A\) 在 \(V\) 上的相应系数范数。
\end{theorem}
\begin{proof}
\(m=0\) 是 \cref{b2:thm:16-interior-h2} 及有限覆盖的结论。归纳假设已给出 \(u\in H_{\mathrm{loc}}^{m+1}\)。对 \(|\beta|=m\)，令 \(w=D^\beta u\in H_{\mathrm{loc}}^1\)。分布的乘积法则给出
\[
 -\operatorname{div}(A\nabla w)
 =D^\beta f+\sum_{0<\gamma\le\beta}\binom\beta\gamma
       \operatorname{div}\bigl((D^\gamma A)\nabla D^{\beta-\gamma}u\bigr).
\]
当 \(\beta=0\) 时和为空，已经处理。对和内每一项，系数 \(D^\gamma A\) 及其一阶导数有界，而 \(\nabla D^{\beta-\gamma}u\) 及其一阶导数的阶数最多为 \(m+1\)，故属于局部 \(L^2\)。乘积因此属于局部 \(H^1\)，其散度属于局部 \(L^2\)。右端第一项也由 \(f\in H^m\) 属于局部 \(L^2\)。将二阶定理应用于 \(w\)，得 \(D^\beta u\in H_{\mathrm{loc}}^2\)，即 \(u\in H_{\mathrm{loc}}^{m+2}\)。

为得到定量界，在 \(U,V\) 之间选取有限个逐层嵌套的开集。每步二阶估计的右端由前一步已经控制的 \(H^{j+1}\) 范数、\(f\) 的 \(H^j\) 范数及系数的 \(W^{j+1,\infty}\) 范数控制。依次代入归纳估计并合并有限个常数，得到所列不等式。乘积法则可以先局部磨光，再在相应 Sobolev 范数中取极限，所以没有预设经典高阶导数。
\end{proof}

对常系数 Laplace 方程，上述论证尤其直接：\(-\Delta D^\beta u=D^\beta f\)。一般变系数方程的额外项只含较低阶的解导数，正是归纳能够闭合的原因；若求导后又在右端出现尚未控制的最高阶导数，就必须先把它保留在主部，不能把它当成已知数据。

\ProseParagraphBreak

若包含低阶项且 \(b,c\) 具有相应有界导数，同样的归纳适用。对非线性方程则须额外检查复合函数的 Sobolev 正则性：从 \(u\in H^2\) 不能在所有维数自动推出 \(F(u)\in H^1\) 或 \(H^2\)。例如含高次幂的项，其导数是 \(F'(u)\nabla u\)，必须先用嵌入定理保证乘积可积。自举是一系列空间之间的合法映射，而不是只按形式导数计数。

\subsection{从弱解到经典解}
\label{b2:subsec:160203}

Sobolev 正则性最终通过嵌入转成经典正则性。第一卷定理 24.4.3 及处理正整数差的命题 24.4.4，在局部截断后给出：若整数 \(s>k+d/2\)，则 \(H_{\mathrm{loc}}^s\) 函数有 \(C^k\) 代表元；更精细地，可取得所有满足 \(k+\alpha<s-d/2\)、\(0<\alpha<1\) 的局部 \(C^{k,\alpha}\) 正则性。严格不等号不能随意替成等号。

\begin{corollary}[光滑数据的内部光滑性]
\label{b2:cor:16-smooth-interior}
设 \(A,f\in C^\infty(\Omega)\)，\(A\) 在每个紧子集附近一致椭圆。若 \(u\in H_{\mathrm{loc}}^1(\Omega)\) 弱解 \(-\operatorname{div}(A\nabla u)=f\)，则 \(u\) 有 \(C^\infty(\Omega)\) 代表元，并逐点满足方程。
\end{corollary}
\begin{proof}
固定 \(U\Subset\Omega\)。光滑系数在稍大紧集上的各阶导数有界，高阶正则性定理给出任意阶的局部 \(H^s\) 正则性。给定 \(k\)，选整数 \(s>k+d/2\)，嵌入得到 \(C^k\) 代表元；各代表元因几乎处处相同且连续而彼此一致，故得到单一 \(C^\infty\) 代表元。该代表元对应的微分表达式连续，与零分布相等，所以处处为零。
\end{proof}

若只知道 \(f\in L^2\)，在 \(d=2,3\) 时 \(H_{\mathrm{loc}}^2\) 已给出一定 Hölder 连续性，却一般不给出连续一阶导数；在 \(d\ge4\) 时甚至不能单凭这一 Sobolev 结论推出连续性。维数越高，要通过平方可积导数达到同一经典阶数，就需要越多次提升。低正则椭圆理论提供了另一条路：即使没有系数导数，也能对实标量方程控制解的振荡；这将在本章最后说明。

\ProseParagraphBreak

“弱解变成经典解”还需要明确所指区域。这里所有结论均为内部结论：边界值不光滑、区域有角点时，解可能在每个内部点都光滑，却在靠近边界处出现不可积的二阶导数。下一节将通过边界坐标和一个扇形反例说明这种差别。

% !TeX root = ../../Book2.tex
\section{边界正则性}
\label{b2:sec:1603}

内部差商可以向任意坐标方向平移，靠近边界时却可能越界。对于零 Dirichlet 条件，平坦边界允许沿切向作差商，而法向二阶导数则由方程补回。曲边界先化成平坦边界；这个坐标变换消耗的正则性，决定了边界应满足的假设。

\subsection{边界平坦化}
\label{b2:subsec:160301}

称有界区域为 \(C^{1,1}\) 区域，是指边界在有限个局部正交坐标中可写成 \(x_d=\psi(x')\)，其中 \(\psi\) 一阶可微且梯度 Lipschitz 连续，区域位于图像的一侧。局部使用
\[x=\Phi(y)=(y',y_d+\psi(y')),\qquad y_d>0\]
把半球或半盒映到区域内部。映射及其逆一阶导数有界，二阶弱导数也有界；因此它们保留局部 \(H^1\)、\(H^2\) 正则性，并给出相应范数的双向控制。

\ProseParagraphBreak

若 \(v(y)=u(\Phi(y))\)、\(\zeta(y)=\varphi(\Phi(y))\)，链式法则和换元把型写成
\[
 \int A(x)\nabla_xu\cdot\nabla_x\varphi\,\mathrm dx
 =\int\widetilde A(y)\nabla_yv\cdot\nabla_y\zeta\,\mathrm dy,
\]
其中
\[\widetilde A=|\det D\Phi|\,D\Phi^{-1}(A\circ\Phi)D\Phi^{-\mathsf T},
 \qquad\widetilde f=(f\circ\Phi)|\det D\Phi|.\]
图像平移的行列式为 \(1\)，保留一般写法便于检查其他坐标。由于 \(D\Phi^{-\mathsf T}\xi\) 的长度与 \(|\xi|\) 双向可比，\(\widetilde A\) 仍一致椭圆；若 \(A\in W^{1,\infty}\)、\(\Phi\in C^{1,1}\)，则 \(\widetilde A\in W^{1,\infty}\)。右端 \(\widetilde f\) 仍属于 \(L^2\)。这正是下一步差商证明需要的三个条件。

\ProseParagraphBreak

只有 Lipschitz 边界时，\(D\Phi\) 一般只有本性有界性，变换后的系数未必具有有界一阶导数。此时切向差商中的系数误差便不再受控。边界光滑性不是为了画图方便，而是通过坐标系数直接进入估计。

\subsection{Dirichlet 问题}
\label{b2:subsec:160302}

先考虑半球 \(B_2^+=B(0,2)\cap\{x_d>0\}\)。设 \(u\in H^1(B_2^+)\) 在平坦部分的迹为零，且弱解 \(-\operatorname{div}(A\nabla u)=f\)。测试函数可以支撑到平坦部分，但须在那里取零迹，并在球面部分附近消失。这样的测试等式由原 Dirichlet 弱形式经局部化得到。

\begin{lemma}[边界切向估计]
\label{b2:lem:16-tangential-estimate}
设 \(A\in W^{1,\infty}(B_2^+)\) 有界且一致椭圆，\(f\in L^2(B_2^+)\)，\(u\) 如上具有零平坦边界迹。则对 \(k<d\) 及 \(j\le d\)，\(\partial_k\partial_j u\in L^2(B_1^+)\)，且
\[\sum_{k<d,j\le d}\lVert\partial_k\partial_j u\rVert_{L^2(B_1^+)}
 \le C\bigl(\lVert f\rVert_{L^2(B_2^+)}+\lVert u\rVert_{H^1(B_2^+)}\bigr).\]
\end{lemma}
\begin{proof}
取在 \(B_1\) 上为 \(1\)、支集在 \(B_{3/2}\) 内的光滑截断 \(\eta\)。对切向方向 \(e_k\)，平移不改变 \(x_d\)，所以 \(D_h^ku\) 的平坦边界迹仍为零。于是 \(-D_{-h}^k(\eta^2D_h^ku)\) 是合法零迹测试函数；它在曲面边界附近为零，没有额外边界项。

在半球上重复内部估计的离散分部积分。所有平移都沿切向，故积分换元不越过平坦边界；截断又避免了球面边界。主项仍有椭圆下界，系数差商仍由 \(\lVert\nabla A\rVert_\infty\) 控制，右端仍由差商的 \(H^1\) 界控制。因此 \(\eta\nabla D_h^ku\) 的 \(L^2\) 范数独立于 \(h\) 有界。差商引理在半球内部测试并取极限，得到所列切向及混合导数及范数界。
\end{proof}

现在只缺 \(\partial_{dd}u\)。在分布意义展开方程，得
\[
 a_{dd}\partial_{dd}u
 =-f-\sum_{(i,j)\ne(d,d)}a_{ij}\partial_{ij}u
       -\sum_{i,j}(\partial_i a_{ij})\partial_j u.
\]
右端各项均属于 \(L^2(B_1^+)\)：所有非纯法向二阶项已由引理控制，系数导数有界，解的一阶导数平方可积。又 \(a_{dd}\ge\lambda>0\)，故除以它得到 \(\partial_{dd}u\in L^2\)。展开与除法均可在分布中合法进行，因为 \(a_{dd}\) 及其倒数是局部 Lipschitz 乘子。至此得到局部边界 \(H^2\) 估计。

\ProseParagraphBreak

这一步没有沿法向把 \(u\) 随意平移，也没有先把 \(u\) 作二阶可微延拓。零迹保证切向测试合法，椭圆性保证纯法向最高阶项的系数可逆，两者承担不同职责。一维时没有切向方向，最后的方程直接恢复二阶导数，恰与一维计算一致。

\subsection{边界 \texorpdfstring{\(H^2\)}{H2} 估计}
\label{b2:subsec:160303}

\begin{theorem}[全局 Dirichlet 二阶正则性]
\label{b2:thm:16-boundary-h2}
设 \(\Omega\subset\mathbb R^d\) 是有界 \(C^{1,1}\) 区域，\(A\in W^{1,\infty}(\Omega)\) 有界且一致椭圆，\(f\in L^2(\Omega)\)。若 \(u\in H_0^1(\Omega)\) 弱解 \(-\operatorname{div}(A\nabla u)=f\)，则 \(u\in H^2(\Omega)\)，且
\[\lVert u\rVert_{H^2(\Omega)}\le C\bigl(\lVert f\rVert_2+\lVert u\rVert_{H^1(\Omega)}\bigr)\le C'\lVert f\rVert_2.\]
常数依赖区域的有限边界坐标、系数上界与 Lipschitz 范数及椭圆下界。
\end{theorem}
\begin{proof}
覆盖 \(\partial\Omega\) 的有限个 \(C^{1,1}\) 坐标把方程变成上一小节的半球模型。坐标变换保留零迹、\(L^2\) 右端及 Lipschitz 椭圆系数；切向估计与恢复纯法向导数给出每个较小边界坐标内的 \(H^2\) 界。余下的紧内部区域由有限个球覆盖，应用内部二阶估计。将有限个局部积分相加，并使用坐标变换的双向范数控制，得到第一项不等式。上一章的能量估计与 Poincaré 不等式控制 \(\lVert u\rVert_{H^1}\le C\lVert f\rVert_2\)，得到第二项。
\end{proof}

对非齐次边界，若给定延拓 \(g\in H^2(\Omega)\)，令 \(w=u-g\)。右端变为 \(f+\operatorname{div}(A\nabla g)\in L^2\)，故
\[\lVert u\rVert_{H^2}\le C\bigl(\lVert f\rVert_2+\lVert g\rVert_{H^2}\bigr).\]
这里比弱存在性所需的 \(g\in H^1\) 多一阶正则性。若只有粗糙边界数据，内部解仍可能光滑，但上述全局二阶结论不再由本定理保证。

\ProseParagraphBreak

对足够光滑的区域和数据，可以在边界坐标中继续切向微分，再用方程恢复高阶法向导数。每次操作都要控制新的系数和坐标导数，因此需要同步提高边界光滑性。全局高阶定理不能仅由内部自举一句带过；本章只完整证明后续最常用的 \(H^2\) 层次。

\subsection{边界光滑性的作用}
\label{b2:subsec:160304}

取平面扇形 \(0<r<1\)、\(0<\theta<\omega\)，其中 \(\pi<\omega<2\pi\)，令 \(\alpha=\pi/\omega\in(1/2,1)\)。函数
\[v(r,\theta)=r^\alpha\sin(\alpha\theta)\]
在扇形内部调和，并在两条直边上为零。它的梯度量级为 \(r^{\alpha-1}\)，故近角点的梯度平方积分由 \(\int_0^\varepsilon r^{2\alpha-1}\,\mathrm dr\) 控制而有限；二阶导数量级为 \(r^{\alpha-2}\)，相应积分 \(\int_0^\varepsilon r^{2\alpha-3}\,\mathrm dr\) 发散。

\ProseParagraphBreak

选径向光滑截断 \(\chi\)，使它在角点附近为 \(1\)，在外圆边界附近为零，令 \(u=\chi v\)。则 \(u\in H_0^1\)，但 \(u\notin H^2\)；源项 \(f=-\Delta u\) 只支撑在离开角点的环带内，因而光滑且平方可积。局部化并在内圆 \(r=\varepsilon\) 上积分，边界余项按 \(\varepsilon^\alpha\) 消失，验证 \(-\Delta u=f\) 确实弱成立。于是即便主部是常系数、源项很好，凹角仍可破坏全局 \(H^2\) 正则性。

\ProseParagraphBreak

这个区域具有 Lipschitz 边界，却在凹角处不属于 \(C^{1,1}\)。例子不意味着所有非 \(C^{1,1}\) 区域都失败，例如某些凸多边形仍有更好的结论；它说明仅凭任意 Lipschitz 边界不能保证本节的全局二阶估计。内部估计在离角点有正距离的每个子域上仍完全成立，两个结论并不冲突。

% !TeX root = ../../Book2.tex
\section{Caccioppoli 不等式与低正则理论初步}
\label{b2:sec:1604}

有界可测系数不足以支持前面的差商估计，却仍保留椭圆性带来的能量控制。对实标量方程，截断与 Sobolev 不等式还能把积分控制逐步变成点态控制。本节完整证明局部有界性，随后明确区分已经完成的能量论证和通向 Harnack、Hölder 理论的深层步骤。

\subsection{Caccioppoli 不等式}
\label{b2:subsec:160401}

以下矩阵 \(A\) 仅假定可测，满足 \(\lVert A\rVert_\infty\le\Lambda\) 和 \(\xi^{\mathsf T}A(x)\xi\ge\lambda|\xi|^2\)。系数不需要一阶导数。

\begin{proposition}[内部能量估计]
\label{b2:prop:16-caccioppoli}
设 \(B_R\Subset\Omega\)，\(u\in H^1(B_R)\) 弱解 \(-\operatorname{div}(A\nabla u)=f\)，\(f\in L^2(B_R)\)。则对 \(0<r<R\) 和任意常数 \(c\in\mathbb R\)，
\[
 \int_{B_r}|\nabla u|^2\le
 \frac{C}{(R-r)^2}\int_{B_R}|u-c|^2
       +C(R-r)^2\int_{B_R}|f|^2,
\]
其中 \(C\) 只依赖 \(d,\lambda,\Lambda\)。
\end{proposition}
\begin{proof}
取 \(\eta\in C_c^\infty(B_R)\)，在 \(B_r\) 上为 \(1\)，满足 \(0\le\eta\le1\)、\(|\nabla\eta|\le C_d/(R-r)\)。函数 \(\eta^2(u-c)\) 是合法测试函数。展开并用椭圆性，得
\[\lambda\int\eta^2|\nabla u|^2
 \le2\Lambda\int\eta|\nabla u||u-c||\nabla\eta|
       +\int|f|\eta^2|u-c|.\]
第一项由 Young 不等式吸收一半左端，余下是 \(C\int|u-c|^2|\nabla\eta|^2\)。对第二项使用带尺度的 Young 不等式，界为 \(C(R-r)^2\int|f|^2+C(R-r)^{-2}\int|u-c|^2\)。整理并利用 \(\eta=1\) 于内球即可。
\end{proof}

这项 \term[caccioppoli-budengshi]{Caccioppoli 不等式}[Caccioppoli inequality]以较大区域内函数的大小控制较小区域内的梯度。取常数 \(c\) 可以去掉与方程无关的整体平移；齐次方程中常取 \(c\) 为球平均值，以突出振荡而不是绝对高度。估计方向与 Poincaré 不等式相反，但需要解满足方程，因此不能对任意 \(H^1\) 函数成立。

\subsection{反向估计}
\label{b2:subsec:160402}

对于齐次解，在 \(B_{2r}\) 上取 \(c=u_{B_{2r}}\)。Caccioppoli 与球上的 Sobolev--Poincaré 不等式结合，给出一个梯度的反向 Hölder 型估计。例如 \(d>2\) 时取 \(q=2d/(d+2)<2\)，有
\[
 \left(\frac1{|B_r|}\int_{B_r}|\nabla u|^2\right)^{1/2}
 \le C\left(\frac1{|B_{2r}|}\int_{B_{2r}}|\nabla u|^q\right)^{1/q}.
\]
这里积分除以球的体积，均取平均值。证明中先用能量估计将左端界为
\[
 \frac Cr\left(\frac1{|B_{2r}|}\int_{B_{2r}}|u-u_{B_{2r}}|^2\right)^{1/2},
\]
再用 Sobolev--Poincaré 将这一项界为右端。二维可选任意固定 \(1<q<2\)，由对应有限指数嵌入及 Hölder 不等式得到同样形式。

\ProseParagraphBreak

通常 Hölder 不等式以较高指数范数控制较低指数范数；这里对解的梯度出现相反方向，代价是扩大球。它暗示梯度具有比任意 \(L^2\) 函数更强的分布控制。由此继续推出高于 \(2\) 的可积指数，需要反向 Hölder 不等式的自改进定理；本节没有证明这一步，因而也不把更高可积性写成已经建立的结论。

\ProseParagraphBreak

局部有界性的迭代使用另一种增强：把能量估计应用于超过某个高度的部分。称 \(u\in H^1\) 为弱次解，是指 \(\int A\nabla u\cdot\nabla\varphi\le0\) 对所有非负紧支集测试函数成立。取 \(w=(u-k)_+\)，经截断逼近可用 \(\eta^2w\) 测试，得到
\[\int\eta^2|\nabla w|^2\le C\int w^2|\nabla\eta|^2.\]
在 \(\{u\le k\}\) 上 \(\nabla w=0\)，在其余部分 \(\nabla w=\nabla u\)，所以与前面的吸收步骤完全相同。高度 \(k\) 可以随迭代改变，这为同时控制高值集合的大小和函数值提供了自由度。

\subsection{局部有界性的动机}
\label{b2:subsec:160403}

\begin{theorem}[弱次解的局部上界]
\label{b2:thm:16-local-boundedness}
设 \(d\ge2\)，\(A\) 在 \(B_R\) 上有界可测且一致椭圆，\(u\in H^1(B_R)\) 是实弱次解。则
\[\operatorname*{ess\,sup}_{B_{R/2}}u_+
 \le C\left(\frac1{|B_R|}\int_{B_R}u_+^2\right)^{1/2},\]
其中 \(C\) 只依赖 \(d,\lambda,\Lambda\)。齐次弱解同时满足相应的 \(|u|\) 局部上界。
\end{theorem}
\begin{proof}
先缩放到 \(R=1\)。对支撑在单位球的 \(H_0^1\) 函数，选 \(\chi=d/(d-2)>1\)（\(d>2\)）或 \(\chi=2\)（\(d=2\)），Sobolev 不等式给出 \(\lVert v\rVert_{2\chi}^2\le C\lVert\nabla v\rVert_2^2\)。固定待选高度 \(K>0\)，令
\[r_n=\tfrac12+2^{-n-1},\qquad k_n=K(1-2^{-n}),\qquad
 Y_n=\int_{B_{r_n}}(u-k_n)_+^2.\]
选截断 \(\eta_n=1\) 于 \(B_{r_{n+1}}\)，支集在 \(B_{r_n}\)，且 \(|\nabla\eta_n|\le C2^n\)。令 \(w_n=(u-k_n)_+\)，截断能量估计连同乘积导数公式给出
\[\lVert\eta_nw_n\rVert_{2\chi}^2
 \le C\lVert\nabla(\eta_nw_n)\rVert_2^2\le C4^nY_n.\]
在集合 \(E_{n+1}=B_{r_{n+1}}\cap\{u>k_{n+1}\}\) 上，\(w_n\ge k_{n+1}-k_n=K2^{-n-1}\)。所以 \(|E_{n+1}|\le4^{n+1}K^{-2}Y_n\)。再用 Hölder 不等式，记 \(\theta=1-1/\chi>0\)，得
\[Y_{n+1}\le\lVert\eta_nw_n\rVert_{2\chi}^2|E_{n+1}|^\theta
 \le Cb^nK^{-2\theta}Y_n^{1+\theta},\]
其中 \(b>1\) 为固定常数。

令 \(Z_n=Y_n/K^2\)，则 \(Z_{n+1}\le Cb^nZ_n^{1+\theta}\)。若
\[Z_0\le C^{-1/\theta}b^{-1/\theta^2},\]
归纳可证 \(Z_n\le C^{-1/\theta}b^{-n/\theta-1/\theta^2}\)：代入递推式后，常数幂及 \(b\) 的指数恰好得到下一项。因此 \(Y_n\to0\)。取 \(K\) 为 \(\lVert u_+\rVert_{L^2(B_1)}\) 的足够大固定倍数即可满足初始小量条件；零范数情形结论显然。因为 \((u-K)_+^2\) 在 \(B_{1/2}\) 上的积分不超过每个 \(Y_n\)，故 \(u\le K\) 几乎处处。缩放回半径 \(R\) 得到平均形式的估计。对于解，\(u\) 与 \(-u\) 都是次解，分别估计正部即可。
\end{proof}

这一证明称为截断迭代。它没有控制系数的导数，也没有先假设解有界；有界性由高值集合的递推衰减产生。迭代中的超线性指数 \(1+\theta>1\) 来自 Sobolev 积分增益，正是小量逐步改善的原因。若只得到线性递推 \(Y_{n+1}\le Cb^nY_n\)，几何增长的常数会使论证失败。

\ProseParagraphBreak

一维可以直接使用 \(H^1\) 函数的绝对连续代表和基本定理获得局部控制，无须上述多维迭代。对于含源项的方程，截断估计会多出源项积分；要保留局部上界，源项需具有与维数相容的可积性。例如一般高维 \(L^2\) 源项未必足够。这里的齐次假设不能只从定理标题中读过而在应用时忘掉。

\subsection{De Giorgi–Nash–Moser 理论简介}
\label{b2:subsec:160404}

局部有界性还没有保证振荡随半径缩小而消失。实标量一致椭圆方程的更强结论是：非负解在较小球内的最大值与最小值可相互比较，任意实解因而具有 Hölder 连续代表。De Giorgi、Nash 和 Moser 的工作从不同方向建立了这一低正则理论。本小节介绍结论及其逻辑联系；其中 Harnack 比较的深层步骤不冒充为上一迭代已经证明。

\begin{theorem}[可测系数的 Harnack 比较，深层结论]
\label{b2:thm:16-measurable-harnack}
设 \(A\) 在 \(B_{2R}\subset\mathbb R^d\) 上为实有界可测矩阵，满足固定的一致椭圆上下界。若 \(u\in H^1(B_{2R})\) 为非负弱解 \(-\operatorname{div}(A\nabla u)=0\)，则
\[\operatorname*{ess\,sup}_{B_R}u\le C\operatorname*{ess\,inf}_{B_R}u,\]
其中 \(C\) 只依赖维数和椭圆上下界。
\end{theorem}
\begin{proofsketch}
上一小节已经完整证明次解的局部上界。余下的核心是超解的正性传播：一定比例的点上具有正下界，可以转成较小区域上的统一正下界。用负幂或对数测试建立这一性质，再结合覆盖分解，可以得到一个小正指数的弱 Harnack 下界；将它与相应指数的局部上界结合，就得到所列比较。这里省略正性传播和覆盖迭代的证明，故这一段只是证明路线。Dongsheng Li 与 Kai Zhang 的论文 \cite[Theorem 1.1 and Section 2]{Li2017Harnack} 给出可测系数的比较定理，并说明局部上界、正性传播与比较之间的关系；系统理论可参阅 \cite[Chapter 8]{Gilbarg2001Elliptic}。本章此前的 \(H^2\) 定理和局部有界性均不依赖此深层结论。
\end{proofsketch}

为了看出比较如何控制振荡，设 \(u\) 是 \(B_{2R}\) 内的任意有界实解，记在外球上的本性上、下界为 \(M,m\)，在内球上的为 \(M_1,m_1\)。函数 \(u-m\) 与 \(M-u\) 均为非负解。对二者应用 Harnack 比较，并将两项不等式相加，得到
\[(M-m)+(M_1-m_1)\le C\bigl((M-m)-(M_1-m_1)\bigr).\]
因此
\[\operatorname*{ess\,osc}_{B_R}u\le\vartheta\operatorname*{ess\,osc}_{B_{2R}}u,
 \qquad\vartheta=\frac{C-1}{C+1}<1.\]
在同心球上迭代，得到振荡按半径的某个正幂衰减。取每点缩小球平均值的极限，振荡衰减保证极限存在；重叠球比较保证这些值满足局部 Hölder 估计，并与原函数几乎处处一致。若 \(C=1\)，解在内球为常数，结论更直接。

\ProseParagraphBreak

这套结论的对象是实标量二阶散度方程。一般椭圆系统缺少同样的序结构，不能直接对向量解作正部截断，也不能把标量 Harnack 定理逐分量套用到耦合系统上。系数虽可跳跃，标量解仍可能 Hölder 连续；其梯度却未必连续，甚至不一定有平方可积的二阶导数。低正则理论改善的是解的振荡，与前面依赖系数导数的高阶 Sobolev 理论相互补充。

% 本章小结（不计入编号）
\begin{chaptersummary}
内部高阶正则性从一个与步长无关的差商估计开始。系数的 Lipschitz 控制保证通量差商的误差有界，而更高系数导数保证自举时产生的右端仍在所需空间。Sobolev 嵌入再把积分意义的导数转成经典代表，其指数门槛依赖维数。边界估计还消耗坐标图的正则性；凹角例子说明内部光滑不自动带来全局 \(H^2\)。

\ProseParagraphBreak

仅有界可测系数时，Caccioppoli 不等式仍然成立。截断迭代把高值集合的能量与测度联系起来，完整给出齐次标量弱解的局部有界性。Harnack 比较进一步控制振荡，但需要另行建立超解的正性传播；从比较到 Hölder 衰减的推论则可直接迭代得到。至此，变分存在性、弱型可解性与正则性三项任务已各有明确的假设和证明。
\end{chaptersummary}

比较本编与后面的估计时，可以固定一个实标量 Dirichlet 问题，再逐项查看增加了什么条件、换来了什么结论。\cref{b2:tab:16-dirichlet-levels}中的各行保留自己的空间和假设，并不是把所有条件一并加到同一个定理上。

\begin{table}[htbp]
\centering
\small
\caption[椭圆 Dirichlet 问题的各层结论]{椭圆 Dirichlet 问题的表示、存在性与正则性}
\label{b2:tab:16-dirichlet-levels}
\begin{tabularx}{\textwidth}{@{}p{0.16\textwidth}XX@{}}
\toprule
所处阶段 & 已得到什么 & 应保留什么条件或限制 \\
\midrule
第十三章 & 球、半空间等模型的 Green 表示与 Poisson 公式。 & 表示一个已有解，与证明一般区域上存在解，是不同任务；核、数据及边界取值须满足对应公式的条件。 \\
第十四章 & 在给定仿射边界类中极小化凸能量。 & 二次能量使用 \(H_0^1\)、合法的边界延拓及 \(H^{-1}\) 源项；非二次能量另核增长、自反性及弱下半连续性。 \\
第十五章 & 强制双线性型的弱解存在、唯一及能量估计。 & 系数有界可测且一致椭圆，源项及测试处于相应对偶空间；不要求系数对称，但不能省去强制性或边界类。 \\
第十六章光滑系数部分 & Lipschitz 系数给内部 \(H^2\)；满足边界条件时进一步给全局 \(H^2\)。 & 内部只考察紧含子域；全局另用有界 \(C^{1,1}\) 区域、\(L^2\) 源项，非齐次边界须有 \(g\in H^2\) 且 \(\operatorname{Tr}u=\operatorname{Tr}g\)。凹角会破坏全局结论。 \\
第十六章及附录 H.1 & 对相应齐次实标量方程作 Harnack 比较并得局部 Hölder 性。 & 比较用于非负弱解和有界可测的一致椭圆系数；不依赖系数可微性，也不给出同样的二阶导数结论。 \\
第十八、二十二章 & 全空间常系数二阶导数的 \(L^p\) 控制，以及局部化、冻结系数的步骤。 & \(1<p<\infty\)；齐次导数估计不自动控制整个非齐次 \(W^{2,p}\) 范数。边界、变系数误差及自举起点仍须另核。 \\
\bottomrule
\end{tabularx}
\end{table}

同一个 \(H^1\) 弱解可以在内部光滑，却因粗糙边界数据或区域凹角而不在全局 \(H^2\)。反过来，对只是可测的系数，Harnack 理论给出的正性与振荡控制，也不是差商高阶理论的简写。\secref{b2:sec:1805}与\secref{b2:sec:2203}将把二阶导数估计放到奇异积分、乘子的框架中；它们增加的是可积指数的分析工具，不撤销已有的边界与系数条件。


\BookExercises

各题均为本书自拟。题目分别改变尺度、系数、数据与边界，检验不同正则性结论的边界；有界解的 Liouville 选做题使用本章明示、附录 H 证明的 Harnack 比较。连续边值题综合变分、弱形式与边界数据，说明内部光滑不等于全局二阶正则；最后的选做题训练边界坐标与通量恢复的具体计算。

% 自拟：差商平移约定与离散能量恒等式。
\Needspace{6\baselineskip}
\begin{exercise}\label{b2:ex:16-discrete-energy}
设 \(d\ge1\)、\(1\le k\le d\)，\(u\in H^1(\mathbb R^d;\mathbb R)\)，\(D_h^ku(x)=[u(x+he_k)-u(x)]/h\)。证明当非零实数 \(h\to0\) 时 \(D_h^ku\to\partial_ku\) 于 \(L^2\)，并证明
\[\int_{\mathbb R^d}u(-D_{-h}^kD_h^ku)=\lVert D_h^ku\rVert_2^2.\]
说明为何差商测试中使用负的 \(D_{-h}^k\)。
\end{exercise}
\begin{solution}
对光滑紧支集函数，线段积分公式和导数的平移连续性给出强收敛。一般 \(u\) 先取 \(H^1\) 光滑逼近，利用差商统一界 \(\lVert D_h^k(u-v)\rVert_2\le\lVert\partial_k(u-v)\rVert_2\)，三角不等式把结论推广。恒等式由离散分部积分得到：\(\int u(-D_{-h}^kw)=\int(D_h^ku)w\)，取 \(w=D_h^ku\) 即可。若漏掉负号，主项就成为负的平方范数，无法使用椭圆正性吸收误差。
\end{solution}

% 自拟：系数Lipschitz但不具二阶正则性的明确反例。
\Needspace{6\baselineskip}
\begin{exercise}\label{b2:ex:16-lipschitz-not-h3}
在 \((-1,1)\) 上令 \(u(x)=\int_0^x(1+|t|)^{-1}\,\mathrm dt\)。证明 \(-((1+|x|)u')'=0\) 弱成立，\(u\in H^2(-1,1)\)，但 \(u\notin H^3(-1,1)\)。
\end{exercise}
\begin{solution}
通量恒为 \(1\)，故与任意测试函数导数的积分为零。\(u'=(1+|x|)^{-1}\) 连续且 Lipschitz，其弱导数在两侧分别为 \((1-x)^{-2}\)（\(x<0\)）与 \(-(1+x)^{-2}\)（\(x>0\)），有界，故 \(u\in H^2\)。\(u''\) 从 \(1\) 跳到 \(-1\)，第三分布导数包含 \(-2\delta_0\)，不能属于 \(L^2\)。这说明右端光滑不足以越过系数正则性的限制。
\end{solution}

% 自拟：全空间常系数二阶估计的Fourier核验。
\Needspace{6\baselineskip}
\begin{exercise}\label{b2:ex:16-fourier-h2}
设 \(f\in L^2(\mathbb R^d)\)。用本卷的 \(2\pi\) Fourier 约定求 \((I-\Delta)u=f\) 的唯一 \(L^2\) 分布解，并证明 \(u\in H^2\)、\(\lVert u\rVert_{H^2}\le C_d\lVert f\rVert_2\)。
\end{exercise}
\begin{solution}
定义 \(\widehat u(\xi)=\widehat f(\xi)/(1+4\pi^2|\xi|^2)\)。该函数属于 \(L^2\)，Plancherel 给出 \(u\)，且乘子恒等式验证分布方程。对 \(|\alpha|\le2\)，函数 \((2\pi i\xi)^\alpha/(1+4\pi^2|\xi|^2)\) 有界，故弱导数 \(D^\alpha u\in L^2\)，并有相应范数界；对有限个指标求和得到估计。若另一 \(L^2\) 解存在，差的 Fourier 变换被严格正的 \(1+4\pi^2|\xi|^2\) 乘后为零，故差为零。
\end{solution}

% 自拟：L2右端在高维不足以保证局部有界。
\Needspace{6\baselineskip}
\begin{exercise}\label{b2:ex:16-l2-unbounded}
在 \(\mathbb R^5\) 中取 \(\chi\in C_c^\infty(B_1)\)，在原点附近为 \(1\)，令 \(u(x)=\chi(x)|x|^{-1/4}\)。证明 \(u\in H_0^1(B_1)\)、\(-\Delta u\in L^2(B_1)\)，但 \(u\) 无界。解释为何这不违反本章局部有界性定理。
\end{exercise}
\begin{solution}
原点附近 \(|u|\) 量级为 \(r^{-1/4}\)，梯度量级为 \(r^{-5/4}\)；乘上五维体积因子 \(r^4\) 后，两者平方均可积。截去一个小球并平滑，误差的 \(H^1\) 范数趋零，故紧支集函数 \(u\) 属于 \(H_0^1\)。由 \(\Delta r^\alpha=\alpha(\alpha+3)r^{\alpha-2}\)，主奇性为 \(r^{-9/4}\)，其平方积分近零为 \(\int r^{-1/2}\,\mathrm dr<\infty\)。小球边界通量及测试函数误差趋零，故无额外点质量。于是 \(-\Delta u\in L^2\)，但 \(u\to\infty\)。局部有界性定理讨论齐次方程，此例右端非零；高维中 \(L^2\) 源项不足以支持同一结论。
\end{solution}

% 自拟：角点指数与Sobolev阈值。
\Needspace{6\baselineskip}
\begin{exercise}\label{b2:ex:16-corner-threshold}
设 \(0<\omega<2\pi\)，\(\alpha=\pi/\omega\)，\(v=r^\alpha\sin(\alpha\theta)\) 定义在平面扇形。对 \(\alpha\ne1\)，求它在角点附近属于 \(W^{2,p}\)、\(1\le p<\infty\) 的准确条件，并单独解释 \(\alpha=1\) 的情形。
\end{exercise}
\begin{solution}
若 \(\alpha\ne1\)，二阶导数的平方和在角向上有非零有界因子，径向量级为 \(r^{\alpha-2}\)。因此二阶导数的 \(p\) 次积分有限恰当且仅当 \(p(\alpha-2)+1>-1\)，即 \(p(2-\alpha)<2\)。此条件也保证较低阶项可积；等号时对数发散。若 \(\alpha\ge2\)，该严格不等式对一切有限 \(p\) 自动成立。\(\alpha=1\) 时 \(v=r\sin\theta=y\) 是线性函数，二阶导数恒零，不能用只看形式幂次的阈值排除它。
\end{solution}

% 自拟：迭代引理的显式小量阈值。
\Needspace{6\baselineskip}
\begin{exercise}\label{b2:ex:16-iteration-lemma}
设非负实数列 \((Y_n)_{n\ge0}\) 满足 \(Y_{n+1}\le Ab^nY_n^{1+\delta}\)，其中 \(A>0\)、\(b>1\)、\(\delta>0\)。证明当 \(Y_0\le A^{-1/\delta}b^{-1/\delta^2}\) 时，\(Y_n\le A^{-1/\delta}b^{-n/\delta-1/\delta^2}\)。说明仅知道 \(Y_0<1\) 为何不总够用。
\end{exercise}
\begin{solution}
初始步正是给定条件。若第 \(n\) 步成立，代入递推式得到 \(A\) 的幂为 \(1-(1+\delta)/\delta=-1/\delta\)，\(b\) 的幂为 \(n-(1+\delta)n/\delta-(1+\delta)/\delta^2=-(n+1)/\delta-1/\delta^2\)，故归纳闭合。特别地序列趋零。小量阈值依赖 \(A,b,\delta\)；若 \(A\) 很大，某个接近 \(1\) 的初值即使小于 \(1\)，第一次递推也可能增大，所以必须通过选择截断高度或归一化满足定量阈值。
\end{solution}

% 自拟：全空间平方可积齐次解的消失。
\Needspace{6\baselineskip}
\begin{exercise}\label{b2:ex:16-l2-liouville}
设 \(d\ge2\)，\(A\in L^\infty(\mathbb R^d;\mathbb R^{d\times d})\)，对几乎处处 \(x\) 和全部 \(\xi\in\mathbb R^d\) 有 \(|A(x)\xi|\le\Lambda|\xi|\)、\(\xi^{\mathsf T}A(x)\xi\ge\lambda|\xi|^2\)，其中 \(0<\lambda\le\Lambda<\infty\) 不随 \(x\) 变化。若实函数 \(u\in H_{\mathrm{loc}}^1(\mathbb R^d)\cap L^2(\mathbb R^d)\) 弱解 \(-\operatorname{div}(A\nabla u)=0\)，不用系数可微性，证明 \(u=0\) 几乎处处。
\end{exercise}
\begin{solution}
对任意中心和半径 \(R\)，局部有界性给出 \(\operatorname*{ess\,sup}_{B_{R/2}}|u|\le C_dR^{-d/2}\lVert u\rVert_{L^2(\mathbb R^d)}\)。固定任意有限球并使 \(R\to\infty\)，右端趋零，因此该有限球内 \(u=0\) 几乎处处。用可数个球覆盖全空间即可。这里全局共同的椭圆常数保证估计中的 \(C_d\) 不随半径增长。
\end{solution}

% 自拟：深层Harnack输入推出有界解的Liouville性质。
\Needspace{6\baselineskip}
\begin{optionalexercise}\label{b2:ex:16-bounded-liouville}
设 \(d\ge2\)，\(A\in L^\infty(\mathbb R^d;\mathbb R^{d\times d})\)，对几乎处处 \(x\) 和全部 \(\xi\in\mathbb R^d\) 有 \(|A(x)\xi|\le\Lambda|\xi|\)、\(\xi^{\mathsf T}A(x)\xi\ge\lambda|\xi|^2\)，其中 \(0<\lambda\le\Lambda<\infty\) 为固定常数。若实函数 \(u\in H_{\mathrm{loc}}^1(\mathbb R^d)\cap L^\infty(\mathbb R^d)\) 弱解 \(-\operatorname{div}(A\nabla u)=0\)，使用本章的深层 Harnack 结论证明 \(u\) 为常数。明确指出本题比平方可积解的 Liouville 题多用了哪项输入。
\end{optionalexercise}
\begin{solution}
由正文的振荡收缩，存在与球无关的 \(0<\vartheta<1\)，使 \(\operatorname*{ess\,osc}_{B_R}u\le\vartheta^n\operatorname*{ess\,osc}_{B_{2^nR}}u\le2\vartheta^n\lVert u\rVert_\infty\)。令 \(n\to\infty\)，每个球内本性振荡为零。重叠球的常数相同，故全空间几乎处处为同一常数；取 Hölder 代表即处处为常数。本题使用 Harnack 导出的振荡收缩，上一题只用已经完整证明的局部上界与全局 \(L^2\) 衰减。
\end{solution}

% 自拟：综合圆周正交展开、凸 Dirichlet 能量、弱形式及内部/边界正则性的区别；光滑域上的障碍来自边界数据。
\Needspace{6\baselineskip}
\begin{exercise}[连续边值与二阶正则性的限制]\label{b2:ex:16-rough-trace-energy}
在单位圆盘 \(D=\{x\in\mathbb R^2:|x|<1\}\) 上取实函数空间。令
\[
 g(\theta)=\sum_{n=1}^\infty n^{-3/2}\cos(n\theta),\qquad
 u(r,\theta)=\sum_{n=1}^\infty n^{-3/2}r^n\cos(n\theta),\quad 0\leq r<1.
\]
在原点按上述各调和多项式的值定义 \(u(0)=0\)。
\begin{enumerate}
\item\label{b2:item:16-trace-energy-class}证明 \(g\) 连续、\(u\in H^1(D)\cap C(\overline D)\)，且 \(u\) 的 Sobolev 迹等于 \(g\)。计算 \(\int_D|\nabla u|^2\)。
\item\label{b2:item:16-trace-energy-minimum}证明 \(u\) 是凸能量 \(E(v)=\frac12\int_D|\nabla v|^2\) 在仿射类 \(u+H_0^1(D)\) 中的唯一极小元，并写出对全部 \(H_0^1(D)\) 测试函数成立的弱方程。
\item\label{b2:item:16-trace-energy-regularity}证明 \(u\) 在圆盘内部光滑，但 \(u\notin H^2(D)\)。解释这为何不与光滑区域上的全局 Dirichlet 二阶正则性矛盾。
\end{enumerate}
\end{exercise}
\begin{solution}
因为 \(\sum_n n^{-3/2}<\infty\)，两级数分别在圆周和闭圆盘上一致绝对收敛；\(r^n\cos(n\theta)\) 是 \(\operatorname{Re}(x_1+\mathrm ix_2)^n\)，故不存在原点的坐标奇性。于是 \(u\) 连续到边界，边界值为 \(g\)。记其前 \(N\) 项为 \(u_N\)。圆周正交性和极坐标积分给出
\[
 \|u_N\|_{L^2(D)}^2
 =\frac\pi2\sum_{n=1}^N\frac{n^{-3}}{n+1},\qquad
 \|\nabla u_N\|_{L^2(D)}^2
 =\pi\sum_{n=1}^N n^{-2}.
\]
将相同计算用于两个部分和的差，得到 \(H^1\) Cauchy 性。其 \(L^2\) 极限与一致极限相同，所以 \(u\in H^1(D)\)，且
\(\int_D|\nabla u|^2=\pi\sum_{n=1}^\infty n^{-2}\)。圆盘是光滑有界区域，\subsecref{b2:subsec:150103}所用的迹算子在 \(H^1\to L^2(\partial D)\) 间连续；\(u_N\to u\) 于 \(H^1\)，而边界部分和一致趋于 \(g\)，因此 \(\operatorname{Tr}u=g\)。该小节的零迹刻画还说明，\(u+H_0^1(D)\) 恰是全部具有这一 Sobolev 迹的 \(H^1\) 函数。

每个 \(u_N\) 都调和，因而对 \(\varphi\in C_c^\infty(D)\)，有 \(\int_D\nabla u_N\cdot\nabla\varphi=0\)。由 \(H^1\) 收敛取极限，再由 \(C_c^\infty(D)\) 在 \(H_0^1(D)\) 中的稠密性，得到
\[
 \int_D\nabla u\cdot\nabla w=0\qquad(w\in H_0^1(D)).
\]
所以对每个允许扰动 \(w\)，
\[
 E(u+w)=E(u)+\int_D\nabla u\cdot\nabla w
                 +\frac12\|\nabla w\|_2^2
 =E(u)+\frac12\|\nabla w\|_2^2.
\]
这既证明极小性，也由 \(H_0^1\) 上的 Poincaré 不等式证明唯一性：取等要求 \(\nabla w=0\)，从而 \(w=0\)。弱方程、固定迹的变分问题和凸性在这里分别负责识别候选者、限定允许扰动与判断唯一极小。

在任意 \(r\leq R<1\) 上，固定阶导数的级数由 \(\sum_n n^M R^{n-M}\) 的尾部控制，故 \(u\) 内部光滑且逐点调和。这也符合\cref{b2:cor:16-smooth-interior}：单位系数和零右端都光滑，且 \(u\) 已是能量弱解。可是对 \(0<R<1\)，仍由正交性，
\[
 \int_{|x|<R}|\partial_{rr}u|^2\,\mathrm dx
 =\frac\pi2\sum_{n=2}^\infty\frac{n-1}{n}R^{2n-2}.
\]
令 \(R\uparrow1\)，非负级数的极限为无穷。若 \(u\in H^2(D)\)，则几乎处处 \(x\ne0\) 有 \(\partial_{rr}u=e_r^{\mathsf T}D^2u\,e_r\)，其中 \(e_r=x/|x|\)，其平方积分应受 Hessian 的平方积分控制，矛盾。圆盘本身光滑，失败来自边界数据。\cref{b2:thm:16-boundary-h2}要求零 Dirichlet 迹，而此处 \(g\ne0\)；非齐次版本\cref{b2:cor:16-nonhomogeneous-boundary-h2}另要求有 \(H^2\) 的边界提升，不能仅凭边值连续就套用。事实上，若这样的提升存在，该推论就会迫使 \(u\in H^2(D)\)，与上述计算矛盾。
\end{solution}

% 自拟：二维图像坐标下计算几何系数，以通量恢复纯法向导数；不承担正文证明依赖。
\Needspace{8\baselineskip}
\begin{optionalexercise}[边界平坦化与法向通量]\label{b2:ex:16-flattening-flux}
设实函数 \(\psi\in C^{1,1}([-2,2])\)，记 \(M=\|\psi'\|_\infty\)，并令
\[
 Q=(-2,2)\times(0,2),\quad Q_1=(-1,1)\times(0,1),\qquad
 F(s,t)=(s,t+\psi(s)),\quad \Omega_\psi=F(Q).
\]
设 \(f\in L^2(\Omega_\psi;\mathbb R)\)，\(u\in H^1(\Omega_\psi;\mathbb R)\) 弱解 \(-\Delta u=f\)，且在下边界图像 \(F((-2,2)\times\{0\})\) 取零迹。置 \(v=u\circ F\)、\(\widetilde f=f\circ F\)。
\begin{enumerate}
\item\label{b2:item:16-flux-coordinate}计算 \(DF\)、\(\det DF\)，从弱形式导出变换系数
\[
 \widetilde A=\begin{pmatrix}1&-\psi'\\-\psi'&1+(\psi')^2\end{pmatrix}.
\]
证明 \((2+M^2)^{-1}|\xi|^2\le\xi^{\mathsf T}\widetilde A\xi\le(2+M^2)|\xi|^2\)，并说明变换后的测试函数在平坦边界及其余边上的要求。
\item\label{b2:item:16-flux-recovery}假设切向差商已给出 \(v_{ss},v_{st}\in L^2(Q_1)\)。不预设 \(v_{tt}\) 存在，使用 \(q=(1+(\psi')^2)v_t\) 先证明 \(\partial_tq\in L^2(Q_1)\)，再求出 \(v_{tt}\) 的公式和 \(L^2\) 上界。
\item\label{b2:item:16-flux-geometry}在以上论证中，\(\psi'\) 有界和 \(\psi''\) 本性有界各用在何处？为何不能只凭边界为 Lipschitz 就沿用同一二阶估计？
\end{enumerate}
\end{optionalexercise}
\begin{solution}
有
\[
 DF=\begin{pmatrix}1&0\\\psi'&1\end{pmatrix},\qquad
 (DF)^{-1}=\begin{pmatrix}1&0\\-\psi'&1\end{pmatrix},\qquad \det DF=1.
\]
对 \(\zeta\in C_c^\infty(Q)\)，令 \(\varphi=\zeta\circ F^{-1}\)。它属于 \(H_0^1(\Omega_\psi)\)，可由密度用于原弱形式。链式法则和换元给出
\[
 \int_Q\widetilde A\nabla v\cdot\nabla\zeta
 =\int_Q\widetilde f\,\zeta,\qquad
 \widetilde A=(DF)^{-1}(DF)^{-\mathsf T}.
\]
计算矩阵乘积便得题中系数。该实对称矩阵行列式为 \(1\)，迹为 \(2+(\psi')^2\)，且二次型为 \((\xi_s-\psi'\xi_t)^2+\xi_t^2\)，故两个特征值为正，上界不超过 \(2+M^2\)，下界由两者乘积为 \(1\) 得到。变换后的 \(v\) 在 \(t=0\) 上取零迹。切向差商所用测试函数可以接触此平坦边界，但必须在那里取零迹，并在左右边和上边附近消失；这样其逆变换仍在原区域的 \(H_0^1\) 中。

记 \(p(s)=\psi'(s)\)。方程的两个通量分量为 \(v_s-pv_t\) 和 \(-pv_s+(1+p^2)v_t\)。已有混合导数的分布交换律给 \(\partial_tv_s=v_{st}\)，所以
\[
 \partial_tq
 =-\widetilde f-\partial_s(v_s-pv_t)+\partial_t(pv_s)
 =-\widetilde f-v_{ss}+p'v_t+2pv_{st}\in L^2(Q_1).
\]
这里只对已有相应方向 \(L^2\) 弱导数的因子使用该方向的乘积法则：\(\partial_sv_t=v_{st}\)、\(\partial_tv_s=v_{st}\) 已知，而 \(v_{tt}\) 尚未得到。因为 \((1+p^2)^{-1}\) 是有界的 \(W^{1,\infty}\) 函数且不依赖 \(t\)，由 \(v_t=(1+p^2)^{-1}q\) 的法向弱微分得到
\[
 v_{tt}=\frac{-\widetilde f-v_{ss}+\psi''v_t+2\psi'v_{st}}{1+(\psi')^2},
\]
以及
\[
 \|v_{tt}\|_{L^2(Q_1)}\le
 \|\widetilde f\|_{L^2(Q_1)}+\|v_{ss}\|_{L^2(Q_1)}
 +\|\psi''\|_\infty\cdot\|v_t\|_{L^2(Q_1)}+2M\|v_{st}\|_{L^2(Q_1)}.
\]
\(\psi'\) 的上界控制坐标的一阶范数和椭圆常数；\(\psi''\) 的上界控制几何系数的一阶导数，也控制最后的 \(\psi''v_t\)。只有 Lipschitz 图像时，\(\psi'\) 虽本性有界，其弱导数未必为有界函数，切向系数差商和上述乘积界便缺少依据。因此斜率受控与斜率变化受控在二阶正则性中承担不同职责。
\end{solution}
